We have introduced the Joint Episode as a mathematically precise and ontologically primitive temporal unit grounded in the three-layer ontology of systemic persistence. Formally, a Joint Episode is defined as a maximal contiguous interval of sustained low antisynchrony $A(k)$ (Layer 2: relational coherence and permeability) that contains at least one significant excursion of the critical-mass metric $M(k)$ (Layer 1: local intensification of Systemic Tau $\tau_s$ combined with RQA trapping). This definition is directly operational: it yields a detection algorithm, the scalar Joint Score $J = D \cdot \overline{M}$ together with its bias-reduced variant $J_{0.7} = D^{0.7} \cdot \overline{M}$, and a clear criterion for episodes that satisfy the joint L1--L2 conditions.

The bias-reduction embodied in $J_{0.7}$ has proven empirically advantageous across the synthetic regimes examined. By attenuating the contribution of raw duration, $J_{0.7}$ reduces the confounding effect of trivially long low-antisynchrony intervals and yields more stable separation between labeled classes. The figure accompanying the formal definition (Figure~\ref{fig:three_layers_je}) renders the relational structure explicit: bottom-up enabling arrows connect local intensification to scale coherence to the possibility of reorganization, while top-down constraints indicate that higher-level organization restricts admissible configurations at lower layers. The Joint Episode thereby appears as the concrete realization of a \textit{kairos}---a structured window of the present in which persistence acquires the relational conditions required for potential transition.

Controlled synthetic validation on coupled Lorenz oscillators demonstrated that episodes preceding detectable structural reorganization can be statistically distinguished from non-transition episodes under the chosen labeling, with the first reliably positive post-episode shift in antisynchrony obtained in this generator. However, even an oracle experiment supplying perfect ground-truth labels for perturbation windows produced only a modest AUC of approximately 0.476 for the best score. This ceiling indicates that, in the low-dimensional modular systems tested, the Joint Score functions primarily as a diagnostic of L1--L2 joint structure rather than as a strong anticipatory marker of Layer-3 reorganization. The ontological reading must therefore be stated with precision: the Joint Episode identifies a necessary relational-temporal unit; it does not yet demonstrate sufficiency for structural change in minimal testbeds.

These results carry direct implications for future research. The framework now supplies both a formal definition and an initial empirical calibration. Subsequent work must test whether richer dynamics, higher-dimensional attractors, continuous-time formulations, or genuine empirical multivariate series reveal stronger coupling between the Joint Episode and actual reorganization events. The open questions listed in Section 8 organize the most pressing theoretical, methodological, and applied directions. In particular, the move from synthetic modular systems to real-world data and the development of continuous, non-thresholded characterizations appear indispensable if the present construct is to fulfill its promise as a general ontological-temporal unit.

In summary, the Joint Episode provides a non-reductive bridge between local persistence, cross-scale coherence, and the possibility of structural transition. Its detection is operational, its interpretation is constrained by current evidence, and its further development constitutes a concrete research program at the intersection of recurrence quantification, dynamical systems theory, and the metaphysics of becoming.